%% kapa: exam=1 round=26 year=2015 period=1 form=A subject=economics subject_ko=경제원론 seq=26 number=66
%% source: Q-Net 감정평가사 기출문제 (cst00309) artlSeq=5138001 file=extracted/1008/5138001/제26회_감정평가사 1차_기출/제26회_감정평가사 1차_기출/1교시_감평_A형.pdf
%% conversion: kordoc 4.24.51 -> tools/gen1.py (auto); review=proofread-sonnet; flags=none
\begin{question}{66}
폐쇄경제의 IS-LM 모형에서 정부는 지출을 증가시키고, 중앙은행은 통화량을 증가시켰다. 이 경우 나타나는 효과로 옳은 것은? (단, IS곡선은 우하향, LM곡선은 우상향한다.)
\begin{choices}
  \item 국민소득은 증가하고, 이자율은 하락한다.
  \item 국민소득은 증가하고, 이자율은 상승한다.
  \item 국민소득은 증가하고, 이자율의 변화 방향은 알 수 없다.
  \item 국민소득은 감소하고, 이자율은 상승한다.
  \item 국민소득은 감소하고, 이자율의 변화 방향은 알 수 없다.
\end{choices}
\end{question}
\answer{3}
